\documentclass[12pt]{article}
\usepackage[utf8]{inputenc}
\usepackage[T1]{fontenc}
\usepackage{mathpazo}
\usepackage{amsmath}
\usepackage[margin=1.25in]{geometry}
\usepackage{titling}
\usepackage[colorlinks=true,linkcolor=black,citecolor=black,urlcolor=black]{hyperref}

\title{Rain on the Parade\\[0.3em]\large Why Wanting Affection Is Not the Same as Craving Oxytocin}
\author{Flyxion \\ \small Independent Researcher}
\date{August 2026}

\begin{document}
\maketitle

\begin{abstract}
Oxytocin has acquired an unusual position in popular accounts of human social life. Affection, physical touch, romantic attachment, parental bonding, loneliness, and even the desire for companionship are frequently explained through a simple neurochemical vocabulary: people seek contact because they are ``craving oxytocin.'' The expression sounds scientific because oxytocin genuinely participates in mammalian social physiology. Yet the inference from participation to motivation is unwarranted. A molecule can mediate, modulate, or accompany a behavior without becoming the object that the organism seeks. This essay examines the conceptual error underlying the idea of an oxytocin craving and the empirical evidence that makes the simple model difficult to sustain. Human social interaction does not produce a uniform endogenous oxytocin response, oxytocin does not exert uniformly prosocial effects, peripheral measurements do not transparently reveal central signaling, and social motivation emerges from interacting neural, physiological, developmental, and environmental systems. The larger problem is explanatory: biochemical description is often mistaken for causal reduction. To say that oxytocin participates in attachment does not mean that attachment is secretly a desire for oxytocin. Rain participates in the growth of a garden. The garden is not therefore trying to obtain rain.
\end{abstract}

\section{The Attractive Explanation}

There is something irresistible about the phrase \emph{the love hormone}. It takes an enormous domain of human experience and gives it a molecule.

People embrace one another. Oxytocin is involved. Parents hold infants. Oxytocin is involved. Romantic partners touch. Oxytocin is involved. Social isolation feels unpleasant, while reunion can feel relieving. Once these observations have been assembled, a tempting compression presents itself:

\begin{quote}
We seek affection because we crave oxytocin.
\end{quote}

The explanation has nearly everything required for successful popular science. It is short. It is mechanistic. It replaces something psychologically complicated with something physically concrete. Most importantly, it contains enough truth to protect the larger claim from immediate suspicion.

Oxytocin is real. Its involvement in reproduction, lactation, maternal behavior, social recognition, and numerous forms of social processing is real. Oxytocin receptors are distributed through neural systems relevant to social behavior. Manipulating oxytocin signaling can alter behavior under some experimental conditions.

None of this establishes that a lonely person is experiencing an oxytocin deficiency, that the subjective desire for touch represents the detection of such a deficiency, or that obtaining social contact functions principally by restoring oxytocin to some preferred quantity.

The popular explanation therefore performs a small but consequential substitution. It begins with the proposition
\[
\text{oxytocin participates in social behavior}
\]
and quietly replaces it with
\[
\text{social behavior is undertaken to obtain oxytocin}.
\]
These are not equivalent claims.

It is worth being precise about where this substitution actually occurs. Few researchers publishing in the peer-reviewed oxytocin literature assert, in so many words, that social behavior is a homeostatic drive toward a fixed oxytocin target. The claim in that strong form belongs almost entirely to popular science writing, magazine explainers, and everyday folk neuroscience, where the phrase ``love hormone'' does most of the argumentative work. The purpose of what follows is not to caricature a research community that has, in fact, produced much of the evidence against the simple model, but to examine why a claim so poorly supported by that same literature nonetheless circulates so persistently in public discourse.

The distance between them is the subject of this essay.

\section{A Mechanism Is Not Necessarily a Goal}

Consider hunger.

The physiological regulation underlying hunger is enormously complicated, but it is at least meaningful to say that an organism possesses regulatory systems responsive to nutritional state. Eating changes variables whose condition participates in producing the motivation to eat. Similar reasoning applies to thirst, thermoregulation, and respiration. These systems provide intuitive models for homeostatic drives because a regulated physiological condition can become displaced and behavior can help restore it.

It is tempting to treat every neurochemical correlate of motivation according to this template.

Suppose affectionate contact changes oxytocin signaling. Suppose further that manipulating the oxytocin system changes some social behaviors. Neither fact demonstrates the existence of an oxytocin homeostat whose deviation is experienced as loneliness and whose restoration is experienced as social satisfaction.

The distinction is between a \emph{mediator} and an \emph{object of regulation}.

A physiological process can be necessary for producing an experience without being the thing that the experience is attempting to obtain. Photoreceptors are necessary for ordinary vision; people looking at sunsets are not therefore seeking photoreceptor activation. Motor neurons participate in walking; the person walking home is not trying to stimulate motor neurons. Dopaminergic systems participate in reinforcement and motivation; this does not establish that every rewarded behavior is fundamentally undertaken to acquire dopamine.

The same caution applies to oxytocin.

A person may want a particular person.

A child may want a parent.

Someone frightened may want reassurance.

Someone lonely may want company.

Someone in love may want physical proximity.

Someone grieving may want a person who can no longer be present.

These desires have objects. Replacing those objects with the neurochemistry involved in realizing the desires does not deepen the description automatically. Sometimes it merely changes the level at which the phenomenon is described and then mistakes that change of level for an explanation.

None of this requires denying that social motivation has homeostatic-\emph{like} features. Isolation reliably activates stress physiology; reunion reliably attenuates it. Separation from attachment figures, in humans and in many other mammals, produces measurable changes in cortisol and autonomic tone, and those changes recede with renewed contact. This pattern is genuinely regulatory. But a regulatory pattern implicates whichever systems happen to instantiate it, and the systems most consistently implicated in the felt urgency of social contact are the stress axis and the endogenous opioid system, not a dedicated oxytocin deficit. A homeostatic feel to social wanting is therefore compatible with everything argued here: it shows that something is being regulated, not that the something is oxytocin, and not that the object of the regulation is the hormone rather than the missing person.

\section{What the Evidence Actually Looks Like}

The simple replenishment story makes an empirical prediction that is stronger than is often acknowledged.

If ordinary social contact is sought substantially because it restores an oxytocin deficit, social interaction should produce a reasonably dependable oxytocin response. The response need not be identical in every individual, but the direction should be sufficiently systematic to support the proposed regulatory interpretation.

The literature is considerably less orderly.

A recent systematic review and meta-analysis provides an especially useful test of this expectation \cite{burenkova2023endogenous}. A systematic review and meta-analysis by Burenkova et al.\ examined sixty-three studies of endogenous oxytocin and human social interaction, with fifty-one studies and 3{,}741 participants entering the quantitative synthesis \cite{burenkova2023endogenous}. In causal designs, social interaction did not produce the expected significant overall hormonal response, either in pre--post comparisons or in comparisons between social-interaction and control conditions. Correlational designs did show a small positive association, but the literature exhibited substantial heterogeneity across paradigms, populations, forms of interaction, sampling procedures, and measurement techniques.

This is not evidence that oxytocin is irrelevant to sociality.

It is evidence against pretending that the relationship is simple.

The distinction matters. If hugging reliably filled an oxytocin tank depleted by insufficient hugging, the conceptual model would at least possess a relatively straightforward physiological candidate mechanism. Instead, increases, decreases, null effects, individual variation, contextual variation, and substantial methodological difficulties coexist in the literature.

The rain does not fall on command.

\section{The ``Love Hormone'' Has Bad Manners}

Another difficulty arises from the word \emph{social} itself.

Social behavior is not synonymous with kindness.

Protecting one's child is social behavior. So is distrusting a stranger. Cooperation is social. Competition is social. Courtship, jealousy, conformity, aggression, reconciliation, recognition, territoriality, favoritism, suspicion, and caregiving all occur within social worlds.

Consequently, even if oxytocin were described accurately as a regulator of social processing, it would not follow that it was a chemical of affection.

Research has repeatedly complicated the early popular picture in which administering oxytocin appeared simply to increase trust, empathy, generosity, or attachment. Effects can depend upon the individual, the relationship between participants, prior experience, perceived threat, group membership, and experimental context. This more conditional picture has motivated accounts in which oxytocin modulates the salience of socially relevant cues rather than acting as a general-purpose promoter of prosocial behavior \cite{shamaytsoory2016social}. Under different conditions, alterations in oxytocin signaling have been associated with behaviors that cannot plausibly be gathered under the ordinary meaning of ``love.''

This is not paradoxical once the original metaphor is abandoned.

A system involved in making social information consequential should not be expected merely to make organisms nicer. Social information sometimes says \emph{approach}. Sometimes it says \emph{protect}. Sometimes it says \emph{remember}. Sometimes it says \emph{avoid}. Sometimes it says that this individual matters more than that individual.

Oxytocin's involvement in these processes is scientifically interesting precisely because the biology is richer than the slogan. Indeed, Leng, Leng, and Ludwig's critical review of the literature questions whether the evidence warrants treating oxytocin as a distinctively ``social neuropeptide'' at all \cite{leng2022social}.

Calling it the love hormone is therefore not merely an innocent simplification. The metaphor encourages a particular causal model: more oxytocin means more love, less oxytocin means less love, and a desire for love can consequently be translated into a desire for oxytocin.

Once the first equation fails, the final translation loses its foundation.

\section{The Measurement Problem}

There is an additional complication hidden by everyday claims about oxytocin levels.

Oxytocin operates both centrally and peripherally. Measuring a concentration in blood, saliva, or another peripheral sample does not amount to placing a meter inside the relevant neural circuitry and reading off ``social bonding.''

Sampling procedures differ. Assays differ. Extraction procedures differ. Timing differs. Oxytocin release itself can be temporally complex. The relationship between peripheral measurements and functionally relevant central signaling is not sufficiently simple to support casual statements about someone's current social neurochemical state. The problem is particularly important because peripheral oxytocin can have multiple physiological sources, while assay and extraction methods produce strikingly different concentration estimates \cite{burenkova2023endogenous}.

Yet popular discourse routinely speaks as though the measurement problem had already been solved.

A person says, ``I need a hug.''

The translation arrives:

\begin{quote}
Your oxytocin is low.
\end{quote}

But usually nobody has measured anything.

More importantly, even measurement would not by itself establish the proposed explanation. Discovering a particular concentration would still require demonstrating that the concentration was causally responsible for the subjective state, that the relevant nervous system was regulating that variable toward a target, and that the sought behavior was organized around correcting its deviation.

A biochemical noun has been substituted for an unobserved mechanism.

The resulting statement sounds more precise while sometimes containing less information.

\section{The Intentional Object}

The conceptual problem can be stated more precisely.

Motivated behavior has an object at the level at which the organism acts.

A thirsty animal seeks water. A frightened infant seeks a caregiver. A lonely person seeks social contact. A lover seeks the beloved.

The biological implementation of those motivations can contain thousands of interacting processes. Neurotransmitters are released. Hormonal states change. Autonomic activity shifts. Memories are retrieved. Predictions are generated. Sensory information is integrated. Motor programs are selected.

There is no contradiction between saying that all of these processes are physical and saying that the organism seeks something other than the molecular events through which seeking becomes possible.

The mistake occurs when causal implementation is confused with intentional content.

Let $B$ denote a behavior, $M$ a molecular process participating in the production of that behavior, and $O$ the object toward which the behavior is directed. From
\[
M \rightarrow B
\]
it does not follow that
\[
B \rightarrow \text{acquisition of } M.
\]
Nor does evidence that $O$ changes $M$ establish that $M$ is the hidden object represented by $O$.

A kiss may alter oxytocin signaling.

It does not follow that the kiss is merely an oxytocin-delivery technology.

The difference is not sentimental. It is logical.

\section{The Dopamine Test}

One way to expose the mistake is to substitute another familiar molecule.

Popular explanations frequently say that people scroll through social media because they are ``craving dopamine.'' Video games provide ``dopamine hits.'' Sugar ``gives dopamine.'' Notifications ``release dopamine.'' The vocabulary has become sufficiently familiar that dopamine sometimes appears to function as a universal currency of wanting.

But dopamine participates in the machinery through which organisms learn, predict, select actions, and pursue rewards. Saying that an organism wants something and that dopaminergic processes participate in wanting it does not establish that dopamine itself is what the organism wants.

The argument would otherwise become circular.

Why did the person pursue the reward?

Because dopamine motivated the pursuit.

Why did dopamine motivate the pursuit?

Because the reward was predicted or valued.

What did the person really want?

Dopamine.

The final answer has explained the behavior by replacing the reward with one component of the machinery through which the reward acquired motivational significance.

The oxytocin story performs the same maneuver with a warmer vocabulary.

Dopamine becomes the molecule of pleasure.

Oxytocin becomes the molecule of love.

Cortisol becomes the molecule of stress.

Serotonin becomes the molecule of happiness.

The nervous system becomes a collection of labeled chemical reservoirs, and ordinary experience becomes the conscious surface of a hidden shopping list.

Real neurobiology is considerably less obliging.

\section{Why the Myth Survives}

The oxytocin story persists partly because it satisfies a cultural preference for explanations that descend scales.

A psychological description appears superficial:

\begin{quote}
She misses her husband.
\end{quote}

A biochemical description appears deeper:

\begin{quote}
Her nervous system is seeking oxytocin.
\end{quote}

But scale is not the same thing as explanatory depth.

Suppose someone waits at an airport because her husband is arriving. Describing ATP metabolism in her leg muscles would take us to a physically smaller scale than describing her expectation of his arrival. It would not explain why she is at the airport.

The appropriate explanatory level depends upon the question.

Biochemistry can tell us extraordinary things about how attachment is instantiated, maintained, modified, and integrated with stress, memory, reward, reproduction, and perception. Psychology can characterize the organization of attachment and desire. Developmental science can investigate how relationships acquire their particular forms. Sociology and anthropology can examine the structures through which intimacy is expressed and understood.

None becomes obsolete merely because molecules exist.

Reduction becomes misleading when discovering a lower-level participant is treated as discovering what a higher-level phenomenon ``really is.''

The oxytocin story is not the first time this mistake has been made with unusual literalness. In the 1960s a purified peptide named scotophobin was reported to transfer a learned fear of darkness from trained rats to untrained ones simply by injection, as though a specific piece of psychological content had been distilled into a molecule and could be physically relocated from one brain to another. The episode, and its subsequent loss of credibility, has been chronicled from the inside by a scientist who took part in the search for such ``memory molecules'' \cite{irwin2006scotophobin}. Where the oxytocin case merely conflates a mediator with an object of desire, the scotophobin case shows what happens when that conflation is pushed to its limit and a specific memory is claimed to be, quite literally, a specific chemical.

\section{A Better Biological Story}

Rejecting the phrase ``oxytocin craving'' does not require returning human sociality to some mysterious realm outside biology.

Quite the opposite.

The biological alternative is richer.

Social motivation arises from organisms whose survival and reproduction have depended upon other organisms for millions of years. Mammalian development is profoundly social. Infants depend upon caregivers. Social information predicts safety, danger, nourishment, competition, mating opportunities, coalition, status, learning, and environmental conditions. Neural systems concerned with reward, stress, memory, perception, endogenous opioids, dopamine, serotonin, vasopressin, oxytocin, sex hormones, autonomic regulation, and numerous other processes interact within this history. The association between oxytocin and love is therefore better treated as a relationship among interacting physiological and behavioral systems than as an identity between a molecule and an emotion \cite{carter2022oxytocin}.

Under such conditions, there is no reason to expect ``wanting another person'' to reduce cleanly to replenishing one chemical.

The organism does not first evolve a need for oxytocin and subsequently discover that other organisms are convenient dispensers.

Oxytocin signaling evolved within organisms already embedded in biological problems involving reproduction, offspring, recognition, proximity, threat, coordination, and social information.

The direction of explanation matters.

Other people are not proxies for the molecule.

The molecule is part of the machinery through which other people matter.

A concrete sketch may help. A better biological account of why a child wants a parent does not begin by naming a hormone. It begins with the ecological fact that infant mammals cannot regulate their own physiology, forage, or defend themselves, so proximity to a caregiver is a survival requirement rather than an optional pleasure. It continues with the observation that a nervous system under this evolutionary pressure will recruit whatever machinery is available, reward circuitry, stress circuitry, olfactory and auditory recognition systems, endogenous opioids, oxytocin and vasopressin signaling, to make caregiver proximity motivationally consequential and caregiver absence aversive. None of those systems needs to be, or plausibly is, the sole carrier of the motivation; each contributes a piece of a distributed, overdetermined mechanism whose function is to keep the infant near the person who keeps it alive. The account is biological through and through. It simply refuses to collapse the whole distributed mechanism into the name of one molecule that happens to be easy to measure and easy to market.

\section{Rain on the Parade}

There is a peculiar disappointment in explanations of this kind being challenged.

The love-hormone story is charming.

Two people embrace. Their brains release oxytocin. They feel closer. Loneliness is a shortage; affection replenishes it. Biology has discovered the chemical substance of intimacy.

It is a beautiful little story.

Unfortunately, beauty is not a method.

Oxytocin remains fascinating without being love in molecular form. Its involvement in social behavior becomes more interesting, not less, when its effects are conditional, distributed, interactive, and sometimes counterintuitive. Human attachment becomes more biologically plausible when understood as the product of interacting systems rather than the output of a single chemical dial.

Most importantly, rejecting the oxytocin-craving metaphor does not make affection less physical.

It makes the physical explanation more demanding.

If someone says, ``I miss her,'' neuroscience does not improve the statement merely by replacing \emph{her} with the name of a peptide.

The person may indeed undergo changes in oxytocin signaling when she returns. Dopaminergic activity may change. Stress physiology may change. Endogenous opioid signaling may change. Heart rate, attention, memory, prediction, and behavior may all change.

Those facts concern how a human organism realizes reunion.

They do not show that reunion was secretly sought for the molecules.

Sometimes a person wants another person.

Biology has to explain how that is possible.

It does not get to explain the person away.

\section*{Conclusion}

The claim that people ``crave oxytocin'' contains a legitimate scientific observation inside an illegitimate inference. Oxytocin participates in important aspects of mammalian social physiology. It can influence social recognition, maternal behavior, attachment-related processes, and responses to socially relevant information. Its signaling changes under some forms of social interaction, and experimental manipulation can alter behavior.

None of these findings establishes an oxytocin deficiency theory of loneliness or an oxytocin replenishment theory of affection.

The distinction is elementary but powerful:
\[
\boxed{
\text{A mechanism through which something is wanted}
\neq
\text{the thing that is wanted}
}
\]

Forgetting that distinction produces an entire genre of folk neuroscience in which dopamine becomes pleasure, serotonin becomes happiness, cortisol becomes stress, and oxytocin becomes love.

The molecules deserve better.

So do the phenomena.

\begin{thebibliography}{99}

\bibitem{burenkova2023endogenous}
Burenkova, O. V., Dolgorukova, T. A., An, I., and colleagues.
``Endogenous Oxytocin and Human Social Interactions: A Systematic Review and Meta-Analysis.''
\textit{Psychological Bulletin} 149, nos.~9--10 (2023): 549--579.
doi:10.1037/bul0000402.

\bibitem{irwin2006scotophobin}
Irwin, L. N.
\textit{Scotophobin: Darkness at the Dawn of the Search for Memory Molecules}.
Lanham, MD: Hamilton Books, 2006.

\bibitem{shamaytsoory2016social}
Shamay-Tsoory, S. G., and Abu-Akel, A.
``The Social Salience Hypothesis of Oxytocin.''
\textit{Biological Psychiatry} 79, no.~3 (2016): 194--202.
doi:10.1016/j.biopsych.2015.07.020.

\bibitem{leng2022social}
Leng, G., Leng, R. I., and Ludwig, M.
``Oxytocin---A Social Peptide? Deconstructing the Evidence.''
\textit{Philosophical Transactions of the Royal Society B: Biological Sciences} 377, no.~1858 (2022): 20210055.
doi:10.1098/rstb.2021.0055.

\bibitem{carter2022oxytocin}
Carter, C. S.
``Oxytocin and Love: Myths, Metaphors and Mysteries.''
\textit{Comprehensive Psychoneuroendocrinology} 9 (2022): 100107.
doi:10.1016/j.cpnec.2021.100107.

\end{thebibliography}

\end{document}
